\documentclass[12pt]{article}

% loads config.sty
\makeatletter
\def\input@path{{../../../../tex/}}
\makeatother
\usepackage{config}
\usepackage{macros}

\title{Comparison with Eftedal et al.\ (2025)}
\date{\today}

\begin{document}

\maketitle

\textit{\textbf{Note:} Parts of this document may be AI-written.}

This note checks the predicted relative correlations of \citet{Eftedal2025PNAS} against the path framework derived in \texttt{relative\_covariance.tex}.
Nothing is rederived here; every expression in the ``ours'' column below is that framework's path-coefficient product evaluated on the pedigree the class names, and each was confirmed against an independent \texttt{pathMgr} model of the same pedigree at assortative-mating equilibrium.

\section{Their notation}

Two parameters appear in their expressions.
$m$ is the \emph{genotypic} partner correlation, which is our $\rho_g$.
A \emph{phenotypic} partner correlation --- our $\rho_y$, written $r$ below --- also appears, but only inside the multiplier that half siblings and half first cousins carry, where it enters as the product $rm$.
Their Table~S8 lists predicted genotypic correlations and their Table~S9 the phenotypic ones, the latter being $h^2$ times the former; we work in genotypic correlations $r_g$ throughout and restore the $h^2$ only for the numerical comparison.

They fit $h^2 = 0.503$ and $m = 0.430$, and fit the half-sibling multiplier freely at $c = 0.972$ with a $95\%$ interval of $[0.864,\, 1.127]$.
It is worth noting at the outset that under primary phenotypic assortment at equilibrium the two partner correlations are not free of each other:
\begin{equation}\label{eq:mr}
    m \;=\; h^2 r ,
\end{equation}
so that $m \leq r$ always, and their fitted values imply $r = m/h^2 = 0.855$.

\section{Result}

Table~\ref{tab:eftedal} sets the two sets of expressions side by side.
Six of the seven classes agree exactly.
The one that does not is step siblings.

\begin{table}[H]
    \begin{floatcard}
    \renewcommand{\arraystretch}{2.2}
    \setlength{\tabcolsep}{5pt}
    \small
    \begin{tabular}{l|c|c|c}
        relation
            & largest $\Nmu$
            & their predicted $r_g$
            & ours \\
        \hline
        full sibling
            & $1$ & $\tfrac{1+m}{2}$ & same \\
        1st cousin
            & $1$ & $\big(\tfrac{1+m}{2}\big)^{3}$ & same \\
        2nd cousin
            & $1$ & $\big(\tfrac{1+m}{2}\big)^{5}$ & same \\
        co-cousin
            & $1$ & $m\big(\tfrac{1+m}{2}\big)^{4}$ & same \\
        half sibling
            & $2$ & $\tfrac{1+2m+rm}{4}$ & same \\
        half 1st cousin
            & $2$ & $\big(\tfrac{1+m}{2}\big)^{2}\tfrac{1+2m+rm}{4}$ & same \\
        \hline
        \textbf{step sibling}
            & $\mathbf{3}$ & $m\big(\tfrac{1+m}{2}\big)^{2}$
            & $m\big(\tfrac{1+r}{2}\big)^{2}$ \\
    \end{tabular}
    \caption{%
        Predicted genotypic correlations for the classes of \citet{Eftedal2025PNAS} Table~S9, with their two half-row multipliers already simplified.
        ``Largest $\Nmu$'' is the number of matings the longest mating chain on the path is crossed in.
        Step siblings are the only class in their table for which that number exceeds two, and the only class on which the two frameworks disagree.%
    }
    \label{tab:eftedal}
    \end{floatcard}
\end{table}

Co-cousins are worth singling out, since they are the class one might expect to be most exposed.
Their definition --- a marriage between the uncle of one proband and the aunt of the other --- places the connecting mating between two individuals whose \emph{own} parental couples lie on different chains.
Every mating chain on that path is therefore a single mating, $r$ never enters, and the two expressions coincide.
The child that the connecting step-couple actually has in their Figure~1 is a common descendant of the two probands' lines and so contributes nothing, which we confirmed directly.

\section{Where the step-sibling expressions part}

Step siblings are the one class in which \emph{both} probands' parental couples sit on a single mating chain.
Writing $A$'s parents as $\{f_A, m_A\}$ and $B$'s as $\{f_B, m_B\}$ with $m_A$ mated to $f_B$, that chain is
\begin{equation*}
    f_A \;\text{---}\; m_A \;\text{---}\; f_B \;\text{---}\; m_B ,
    \qquad \Nmu(f_A, m_B) = 3 ,
\end{equation*}
and the class weight is the average of the standardized chain covariance $c$ over the four cross pairs of parents, at mating distances $1, 2, 2$ and $3$.
The disagreement is entirely in how $c$ is evaluated at distances beyond one.
Their expression expands to
\begin{equation*}
    4\,m\left(\frac{1+m}{2}\right)^{2}
    \;=\; m \;+\; 2m^{2} \;+\; m^{3}
    \qquad\Longrightarrow\qquad
    c(1) = m,\quad c(2) = m^{2},\quad c(3) = m^{3},
\end{equation*}
that is, every mating step priced at the genotypic partner correlation.
The chain law is instead $c(n) = m\,r^{\,n-1}$: the first step crosses from a breeding value to a phenotype and costs $m$, and every step after that propagates at the phenotypic level and costs $r$.

This is also where their table stops being self-consistent.
Their half-sibling row uses $4w = 1 + 2m + rm$, which prices a two-step chain crossing at $c(2) = rm$ --- the correct value, imported with Nagylaki's multiplier \citep{Nagylaki1978}.
Their step-sibling row prices the same two-step crossing at $c(2) = m^{2}$.
The two rows cannot both hold unless $r = m$, which by \eqref{eq:mr} would require $h^2 = 1$.

\section{How large the difference is}

The two step-sibling expressions differ by
\begin{equation}\label{eq:gap}
    \frac{\text{ours}}{\text{theirs}} \;=\; \left(\frac{1+r}{1+m}\right)^{\!2} ,
\end{equation}
which is $1$ when $r = m$ and grows as the two partner correlations separate.
At their fitted $h^2 = 0.503$ and $m = 0.430$, equation~\eqref{eq:mr} gives $r = 0.855$ and a ratio of $1.68$.
In phenotypic correlations, against an observed step-sibling correlation of $0.127$ with interval $[0.108,\, 0.147]$:

\begin{table}[H]
    \begin{floatcard}
    \renewcommand{\arraystretch}{1.5}
    \small
    \begin{tabular}{l|ccc}
        $r$ & implied $c$ & our predicted correlation & ratio to theirs \\
        \hline
        $0.430\ (=m)$    & $1.000$ & $0.111$ & $1.00$ \\
        $0.600$          & $1.036$ & $0.138$ & $1.25$ \\
        $0.855\ (=m/h^2)$& $1.089$ & $0.186$ & $1.68$ \\
    \end{tabular}
    \caption{%
        Predicted step-sibling phenotypic correlations at $h^2 = 0.503$, $m = 0.430$, as the phenotypic partner correlation is varied.
        The first row is their published prediction of $0.111$.
        The implied $c$ column is the value their freely fitted half-sibling multiplier would take at that $r$; all three lie inside their interval $[0.864,\,1.127]$.%
    }
    \label{tab:eftedal-numbers}
    \end{floatcard}
\end{table}

\section{What this does and does not change}

Three things bound the practical importance of the correction.

First, it reaches one class only.
Co-cousins are unaffected, and they are both the larger in-law sample --- on the order of $800{,}000$ pairs against roughly $10{,}000$ step-sibling pairs --- and the class carrying their conclusion about relatives-in-law.
Their co-cousin prediction of $0.057$ against an observed $0.075$ stands exactly as they report it.

Second, its size turns on $r/m$, which their data barely constrain.
Their interval on $c$ spans the range in Table~\ref{tab:eftedal-numbers} entirely, and they say as much themselves.

Third, and least comfortably, the correction does not improve the fit.
At the equilibrium value $r = m/h^2$ the corrected prediction of $0.186$ overshoots the observed $0.127$ and falls outside its interval, where their expression undershoots it only slightly.
The right reading of that is probably not that the corrected expression is wrong but that the equilibrium constraint \eqref{eq:mr} is: a phenotypic partner correlation of $0.855$ for school performance is far above any plausible spousal correlation, and their own point estimate $c = 0.972$ implies $r = 0.297 < m$, which \eqref{eq:mr} forbids outright.
Both readings say the same thing --- that their fitted $m$ is larger than primary phenotypic assortment at equilibrium can generate --- which is the direction their paper argues for on other grounds.

It is worth noting what the corrected expression does to the shape of the residuals.
Under their version, the Fisherian model undershoots both in-law classes, which admits a single shared-environment reading.
Under the corrected one it overshoots step siblings while still undershooting co-cousins, and that sign flip within the same category is harder to attribute to one mechanism.

\bibliographystyle{unsrtnat}
\bibliography{../../../../bib/references}

\end{document}
